
\documentclass[a4paper]{article}
\usepackage{graphicx}
\graphicspath{ {C:\Users\AJ\Desktop} }
\begin{document}

Perdidos en la ignorancia 

La idea es una especie de Escape Room. Los jugadores tendrán que ir pasando pruebas por cada habitación que se encuentren (niveles). Al pasar la prueba, podrán acceder a la siguiente pantalla. 
El ranking va por tiempo y nivel conseguido. En las habitaciones puede haber aleatoriamente objetivos “bonus” que al conseguirlos permitirían al jugador saltarse niveles o pistas, reducir el tiempo final de la prueba, etc.
Habrá pruebas contrarreloj aleatoriamente que obligarán al jugador a pasarlas en el tiempo establecido. En caso de no conseguirlo, la partida terminará y deberá volver a empezar.
El juego es individual. No puedes morir, a menos que te encuentres en una sala contrarreloj. Las puntuaciones serán individuales, pero se almacenarán de forma online, es decir, todo el mundo podrá ver los resultados de los demás.  El progreso de los usuarios no se perderá, es decir, la próxima vez que vuelvan podrán continuar por donde lo dejaron. En el juego existirá el rol de administrador que permitirá gestionar los niveles y administrar las cuentas, con privilegios extras.

\begin{figure}[th]
  \centering
\includegraphics[width=0.7\textwidth]{logo.png}
   \caption{Logo del juego}
  \label{fig:gamelogo}
  \end{figure}
  \end{document}
